% META: {"id":"1NSI-TC-EX-016","chapitre":"1NSI-TYPES-CONSTRUITS","type_objet":"exercice","capacites":["1NSI-TYPES-CONSTRUITS-C2"],"parcours":1,"competences":["programmer"],"duree_min":8,"mode_creation":"ex_nihilo","sources_inspiration":[],"status":"approved","origin":"needs_review","status_history":["needs_review","audited","accepted","approved"],"release_acceptance":"RELEASE_OWNER_BATCH_ACCEPTANCE","acceptance_closure_digest":"sha256:9b3ccf9a81c5520fbb7e03b7b2d3e7bf2057a02834b6d908bf3be5f49f3b553f","acceptance_receipt_digest":"sha256:4fad86f0282e0fa4e0419cca1ad6379ae7af5a280c04a4a90880f90a1c044242"}
% BEGIN-VERIFY
% def moyenne(tab):
%     somme = 0
%     for val in tab:
%         somme = somme + val
%     return somme / len(tab)
% assert moyenne([10, 20, 30]) == 20.0
% assert moyenne([15]) == 15.0
% assert abs(moyenne([12, 8, 15, 3, 20]) - 11.6) < 1e-9
% END-VERIFY
\begin{exercice}{1NSI-TC-EX-016}{1}{8}
Écrire une fonction \lstinline{moyenne(tab)} qui prend un tableau non vide d'entiers
et renvoie la moyenne de ses éléments.

\textbf{Exemples :}
\begin{console}
>>> moyenne([10, 20, 30])
20.0
>>> moyenne([12, 8, 15, 3, 20])
11.6
\end{console}
\end{exercice}
